% Options for packages loaded elsewhere
\PassOptionsToPackage{unicode}{hyperref}
\PassOptionsToPackage{hyphens}{url}
\documentclass[
  11pt,
  letterpaper,
]{article}
\usepackage{xcolor}
\usepackage[margin=1in]{geometry}
\usepackage{amsmath,amssymb}
\setcounter{secnumdepth}{-\maxdimen} % remove section numbering
\usepackage{iftex}
\ifPDFTeX
  \usepackage[T1]{fontenc}
  \usepackage[utf8]{inputenc}
  \usepackage{textcomp} % provide euro and other symbols
\else % if luatex or xetex
  \usepackage{unicode-math} % this also loads fontspec
  \defaultfontfeatures{Scale=MatchLowercase}
  \defaultfontfeatures[\rmfamily]{Ligatures=TeX,Scale=1}
\fi
\usepackage{lmodern}
\ifPDFTeX\else
  % xetex/luatex font selection
\fi
% Use upquote if available, for straight quotes in verbatim environments
\IfFileExists{upquote.sty}{\usepackage{upquote}}{}
\IfFileExists{microtype.sty}{% use microtype if available
  \usepackage[]{microtype}
  \UseMicrotypeSet[protrusion]{basicmath} % disable protrusion for tt fonts
}{}
\makeatletter
\@ifundefined{KOMAClassName}{% if non-KOMA class
  \IfFileExists{parskip.sty}{%
    \usepackage{parskip}
  }{% else
    \setlength{\parindent}{0pt}
    \setlength{\parskip}{6pt plus 2pt minus 1pt}}
}{% if KOMA class
  \KOMAoptions{parskip=half}}
\makeatother
\usepackage{longtable,booktabs,array}
\usepackage{calc} % for calculating minipage widths
% Correct order of tables after \paragraph or \subparagraph
\usepackage{etoolbox}
\makeatletter
\patchcmd\longtable{\par}{\if@noskipsec\mbox{}\fi\par}{}{}
\makeatother
% Allow footnotes in longtable head/foot
\IfFileExists{footnotehyper.sty}{\usepackage{footnotehyper}}{\usepackage{footnote}}
\makesavenoteenv{longtable}
\setlength{\emergencystretch}{3em} % prevent overfull lines
\providecommand{\tightlist}{%
  \setlength{\itemsep}{0pt}\setlength{\parskip}{0pt}}
\usepackage{bookmark}
\IfFileExists{xurl.sty}{\usepackage{xurl}}{} % add URL line breaks if available
\urlstyle{same}
\hypersetup{
  hidelinks,
  pdfcreator={LaTeX via pandoc}}

\author{}
\date{}

\begin{document}

\section{The Tholonic Model as a Candidate Alternative to the Standard
Model: Structural Derivations and Parameter
Reduction}\label{the-tholonic-model-as-a-candidate-alternative-to-the-standard-model-structural-derivations-and-parameter-reduction}

\textbf{Author:} Jeffrey W. Milton, Clarity Coalition

\textbf{Date:} 2026

\textbf{Keywords:} tholonic model, Standard Model, N-D-C triad, fine
structure constant, fermion generations, color charge, quark
confinement, Koide formula, parameter reduction, unified field theory,
falsifiability

\textbf{Provisional arXiv subjects:} hep-ph; hep-th (secondary: math-ph;
physics.hist-ph)

\begin{center}\rule{0.5\linewidth}{0.5pt}\end{center}

\subsection{Abstract}\label{abstract}

The Standard Model of particle physics is the most precisely tested
physical theory in history. It is also explicitly incomplete: it
requires 17 independently postulated quantum fields and 19 free
parameters inserted from measurement rather than derived from principle.
A candidate replacement theory must satisfy a specific criterion: it
must reduce the axiom count and derive at least some of the free
parameters from structural geometry. This paper evaluates the tholonic
model against that criterion. The tholonic framework proposes a single
generative structure, the N-D-C triad (Negotiation, Definition,
Contribution), as the irreducible unit of all stable, self-sustaining
organization. Eight prior papers in this series establish the
mathematical grounding of this framework, prove the structural necessity
of exactly three roles (Paper 3), derive five classical mathematical
constants from the triadic recursion (Paper 1), and show that the atom's
measurable properties, including quark occupancy structure, color
charge, the virial theorem, tetrahedral orbital geometry, and the Koide
formula's unexplained \(2/3\) lepton mass ratio, are structural
predictions of the tholonic framework rather than independent empirical
accidents (Paper 8). This paper synthesizes those results into a direct
comparison with the Standard Model. The tholonic framework addresses
five of the Standard Model's unexplained structural facts (three fermion
generations, three color charges, three spatial dimensions, quark
confinement, and the Koide \(2/3\)) through structural arguments that
require no new free parameters. A systematic search over
tholonic-constant expressions identifies
\(4 \cdot e^5 \cdot (\ln 2)^4 = 137.035865\) as the leading numerical
candidate for \(1/\alpha\), with an error of \(-0.98\) ppm and all
exponents being fixed thologram structural values, not fitted
parameters. A second open derivation targets the proton/electron mass
ratio via the known approximation \(6\pi^5 \approx 1836.12\) and its
connection to the tholonic N-weight. The paper includes a systematic
parameter comparison table, explicit falsifiability conditions, and a
prioritized roadmap for the derivations that would, if completed,
constitute the strongest available confirmation of the framework. The
tholonic model is speculative; all claims should be understood as
structurally motivated hypotheses requiring independent verification.

\begin{center}\rule{0.5\linewidth}{0.5pt}\end{center}

\subsection{1. Introduction}\label{introduction}

\subsubsection{1.1 The Standard Model's
achievements}\label{the-standard-models-achievements}

The Standard Model of particle physics is among the most successful
scientific theories ever constructed. It correctly predicts the
anomalous magnetic moment of the electron to eleven significant figures,
the existence and mass of the W and Z bosons, the top quark, and the
Higgs boson before any of them were directly measured. Its quantum field
theoretic framework has been validated across energies ranging from
atomic-scale phenomena (\(\sim 1\) eV) to the highest-energy collider
experiments (\(\sim 10^4\) GeV). No confirmed experimental deviation
from Standard Model predictions has been established despite decades of
intensive search.

This paper does not dispute any of those achievements. The tholonic
model does not predict that any currently confirmed Standard Model
result is wrong. Its claim is structural and prior: that the Standard
Model, as presently constituted, is not a complete explanation of the
physical world, and that the tholonic framework addresses some of its
structural gaps in a way that reduces rather than multiplies explanatory
posits.

\subsubsection{1.2 The Standard Model's structural
limitations}\label{the-standard-models-structural-limitations}

The Standard Model's incompleteness is not a matter of experimental
dispute. It is openly acknowledged by its practitioners and constitutes
one of the primary motivations for theoretical work in fundamental
physics. The specific deficits relevant to this paper are:

\textbf{The parameter problem.} The Standard Model requires 19 free
parameters (26 including the full neutrino sector). None of these is
derived from the theory's structure. Each is measured and inserted. A
table of the 19 parameters is given in Section 5. The electromagnetic
coupling constant \(\alpha \approx 1/137.036\), the ratio of the proton
mass to the electron mass (\(\approx 1836.15\)), and the Koide formula's
ratio \(Q = 2/3\) for lepton masses are among the most prominent
examples. Richard Feynman described \(\alpha\) as ``one of the greatest
damn mysteries of physics.'' No Standard Model calculation produces its
value; it is taken from experiment.

\textbf{The field proliferation problem.} The Standard Model requires 17
independently postulated quantum fields: one for each of the 6 quarks, 3
charged leptons, 3 neutrinos, the photon, the W\(^+\), W\(^-\), Z\(^0\),
gluon, and Higgs boson. These are not derived from each other. Each is a
separate axiom. A theory that replaced the 17 axioms with one and
derived the particle spectrum from the geometry of that one field would,
by any standard measure of explanatory economy, be a better theory.

\textbf{The structural silence problem.} The Standard Model has no
answer to three of the most basic structural questions about the
physical world: Why are there exactly three fermion generations? Why are
there exactly three color charges? Why are there exactly three spatial
dimensions? These are observed facts that the Standard Model
accommodates without explaining.

\textbf{The confinement problem.} The Standard Model derives quark
confinement from the running of the strong coupling constant in QCD.
This is correct and well-tested. But it does not explain why the
confinement principle generalizes: any partial combination of less than
three quarks is either transient or confined. The QCD explanation is
force-specific; the generalization across scales requires additional
argument.

\subsubsection{1.3 The criterion for a better
theory}\label{the-criterion-for-a-better-theory}

Following the standard of explanatory economy articulated by Einstein in
his pursuit of a unified field theory, a candidate replacement or
successor to the Standard Model must satisfy at least one of:

\begin{enumerate}
\def\labelenumi{\arabic{enumi}.}
\tightlist
\item
  Reduce the number of independently postulated fields (axiom
  reduction).
\item
  Derive at least one currently free parameter from structural geometry
  (quantitative derivation).
\item
  Explain at least one currently unexplained structural fact (structural
  prediction).
\end{enumerate}

This paper evaluates the tholonic model against all three criteria. The
claim is not that the tholonic model currently satisfies all three
comprehensively. The claim is that it satisfies criterion 3 on five
counts, is within 1 ppm on criterion 2 for the fine structure constant,
and provides a principled path toward criterion 1.

\begin{center}\rule{0.5\linewidth}{0.5pt}\end{center}

\subsection{2. The Tholonic Framework: Summary of Prior
Results}\label{the-tholonic-framework-summary-of-prior-results}

\subsubsection{2.1 The core claim}\label{the-core-claim}

The tholonic model's foundational claim is that all stable,
self-sustaining structures in nature are instances of a single triadic
pattern called the N-D-C triad: Negotiation (the emergent equilibrium),
Definition (the bounding/constraining role), and Contribution (the
integrating/accumulating role). This is a structural claim, not a
metaphysical one. It asserts that any entity that persists in time does
so because it simultaneously satisfies three functional roles: something
that limits its form (D), something that flows and integrates across
that form (C), and a stable balance state that emerges when D and C are
in equilibrium (N).

Paper 3 of this series proves formally that three roles are the minimum
for non-trivial convergent recursion. A system with fewer than three
roles is either trivially stable (constant, no dynamics) or divergent
(no stable equilibrium). The proof is constructive and does not depend
on any physical assumption.

\subsubsection{2.2 N-D-C role assignments and the
thologram}\label{n-d-c-role-assignments-and-the-thologram}

The thologram is the geometric realization of the N-D-C triad at
successive levels of recursion. It is a triangular structure with six
labeled points: three outer vertices carrying the N, D, C roles and
three inner midpoints carrying complementary roles. Three distinct
numerical patterns arise from this structure:

\textbf{Binary vertex values (1, 2, 4):} The outer vertices carry
\(N = 2^0 = 1\), \(D = 2^1 = 2\), \(C = 2^2 = 4\). These reflect the
roles' positions in a binary state space.

\textbf{Axis multipliers (2, 3, 5):} Summing vertex values along each
directed axis yields \(14 = 7 \times 2\) (N axis), \(21 = 7 \times 3\)
(C axis), \(35 = 7 \times 5\) (D axis). The axis multipliers are the
tholographic seeds for the \(\pi/4\) branch derived in Paper 1.

\textbf{Role weights (2-3-6):} The structural role weights are D = 2, C
= 3, N = 6, with D:C weight ratio \(= 2/3\).

These three patterns are not mutually exclusive. A given physical
correspondence may be illuminated differently by each, and the
appropriate pattern for a given context is determined by which aspect of
the tholonic structure is being examined.

\subsubsection{2.3 Results from the paper
series}\label{results-from-the-paper-series}

The eight prior papers establish the following results on which this
paper builds:

\begin{itemize}
\tightlist
\item
  \textbf{Paper 1:} Five classical constants (\(\pi/4\), \(\phi\),
  \(e\), \(\sqrt{2}\), \(\ln 2\)) emerge as limits of the tholonic
  recursive system from its own internal geometry, without separate
  derivations for each.
\item
  \textbf{Paper 3:} Three functional roles are the proved minimum for
  non-trivial convergent recursion. Any fewer roles produces either a
  constant system or a divergent one.
\item
  \textbf{Paper 6:} The integers 1, 2, and 3 have qualitatively distinct
  structural roles in minimal recursive systems, providing a formal
  basis for the N-D-C assignment.
\item
  \textbf{Paper 7:} The N-D-C framework has been independently
  rediscovered by Cambridge Semantics through two decades of production
  engineering of large-scale knowledge graphs, providing an empirical
  convergence that is not constructed from tholonic premises.
\item
  \textbf{Paper 8:} The atom is measurably a tholon. Six specific
  correspondences are established, graded by strength, and compared with
  the Standard Model's explanatory coverage.
\end{itemize}

\begin{center}\rule{0.5\linewidth}{0.5pt}\end{center}

\subsection{3. Structural Axiom
Reduction}\label{structural-axiom-reduction}

\subsubsection{3.1 The three-role minimality
theorem}\label{the-three-role-minimality-theorem}

Paper 3's central result is that the N-D-C triad is not an arbitrary
choice of three roles but the unique minimum structure for convergent
recursive self-organization. The theorem states: any recursive system
that achieves non-trivial stable convergence requires at least three
distinct functional roles in the N-D-C pattern. Systems with fewer roles
either do not converge (no D to bound C) or converge trivially to a
constant (no C to integrate). Systems with more than three roles are
either decomposable into N-D-C triads or unstable due to
over-constraint.

This result has direct consequences for how many of the Standard Model's
separately postulated structures are actually independent. If the three
fermion generations, three quark colors, and three-quark composition of
hadrons are all consequences of the same three-role minimality
requirement applied at different scales, then they are not three
independent axioms. They are one axiom, instantiated three times.

\subsubsection{3.2 Three color charges as a tholonic structural
prediction}\label{three-color-charges-as-a-tholonic-structural-prediction}

Quantum chromodynamics requires exactly three color charges (red, green,
blue) as a SU(3) gauge symmetry. The Standard Model postulates SU(3) as
the gauge group of the strong force. It does not derive why the gauge
group is SU(3) rather than SU(2) or SU(4).

The tholonic argument for three colors is structural: a stable quark
configuration requires three N-D-C roles, one at each vertex of the
tholographic trigram. Each vertex carries one color. Color neutrality,
the observed requirement that stable hadrons carry zero net color
charge, is the tholonic requirement that every stable N-D-C
configuration must be balanced across all three roles. An unbalanced
configuration (one or two colors dominant) is the tholonic equivalent of
an unstable partial tholon.

The tholonic model does not replace the QCD derivation of SU(3); it
provides a prior structural reason for why a three-generator gauge
symmetry is the correct framework. The convergence of SU(3) from QCD and
three-role minimality from the tholonic model is itself a form of
consilience.

\subsubsection{3.3 Three fermion generations as a structural
prediction}\label{three-fermion-generations-as-a-structural-prediction}

The Standard Model accommodates three generations of fermions (electron
family, muon family, tau family; up/down family, charm/strange family,
top/bottom family) but provides no explanation for why there are exactly
three. No fourth-generation charged fermion has been found; precision
electroweak measurements at LEP constrain the number of light neutrino
species to \(N_\nu = 2.9840 \pm 0.0082\), consistent with exactly three.

The tholonic prediction follows directly from Paper 3's minimality
theorem: if the three-role requirement applies at the generation level
of the fermion spectrum (each generation playing one of the N, D, C
roles in the fermion tholon), then exactly three generations is a
structural necessity, not a numerical coincidence. The N generation
(first: electron, up/down) is the lightest and most stable; the D
generation (second: muon, charm/strange) is the bounding/constraining
generation; the C generation (third: tau, top/bottom) is the heaviest,
most unstable, most integrating generation.

The asymmetry in stability, with lighter generations more stable and
heavier ones unstable on cosmological timescales, is consistent with the
tholonic stability gradient: the N role (first generation) is the most
stable equilibrium state, D (second) is stable but transient under weak
decay, and C (third) is the least stable, most contributing role,
characteristically unstable and decaying to the first generation under
normal conditions.

A fourth generation is not prohibited by experiment at high mass, but is
constrained by the LEP measurement cited above for light neutrinos. The
tholonic model's structural prediction is that no fourth generation
exists at any mass, because the three-role structure is complete. This
is falsifiable by discovery of a fourth-generation fermion.

\subsubsection{3.4 Quark confinement as incomplete-tholon
instability}\label{quark-confinement-as-incomplete-tholon-instability}

QCD derives quark confinement through the running of the strong coupling
constant \(\alpha_s\), which increases with distance (infrared slavery).
As two quarks are pulled apart, the energy in the gluon field between
them grows linearly with distance until it exceeds the energy required
to create a new quark-antiquark pair from the vacuum. Confinement is the
result: free quarks are never observed.

The tholonic structural argument provides an independent reason for the
same fact. A single quark occupies one N-D-C role in the tholographic
trigram. It is, by definition, an incomplete tholon: it has one role but
lacks the two remaining roles required for a stable N-D-C configuration.
An incomplete tholon is structurally unstable. It cannot persist as an
isolated entity; it must either complete its structure (combine with
other quarks to form a hadron) or not exist as a stable isolated object.

The tholonic and QCD explanations are complementary rather than
competing. The QCD explanation specifies the mechanism (running
coupling, string formation, pair creation). The tholonic explanation
identifies the structural reason that generalizes beyond the specific
force: any partial tholon is unstable, regardless of the specific
interaction governing it. This generalization predicts that
confinement-like behavior should appear in any system where the
underlying structure is tholonic and partial configurations are
accessible. The quark system is the clearest known instance.

\subsubsection{3.5 Three spatial dimensions as a tholonic
constraint}\label{three-spatial-dimensions-as-a-tholonic-constraint}

The Standard Model is formulated in 3+1 spacetime but provides no
explanation for why physical space has exactly three dimensions. The
tholonic model offers a structural argument: the tetrahedron is the
minimum structure that encloses a volume and the minimum
three-dimensional closure of the N-D-C trigram. In fewer than three
spatial dimensions, no volume can be enclosed and the tholographic
recursion cannot achieve three-dimensional self-similarity. In more than
three spatial dimensions, the tetrahedron is no longer the minimal
enclosing structure and the tholonic three-role closure is not uniquely
defined.

If the tholonic framework formally proves that stable recursive N-D-C
structure requires exactly three spatial dimensions for volumetric
closure, this constitutes a structural derivation of a fact the Standard
Model cannot address. This argument is currently informal; a formal
proof is identified as an open derivation in Section 7.

\begin{center}\rule{0.5\linewidth}{0.5pt}\end{center}

\subsection{4. Quantitative Parameter
Derivations}\label{quantitative-parameter-derivations}

\subsubsection{4.1 The fine structure constant: leading
candidate}\label{the-fine-structure-constant-leading-candidate}

The fine structure constant \(\alpha\) governs the strength of the
electromagnetic interaction. Its measured value (CODATA 2022) is:

\[\frac{1}{\alpha} = 137.035999206(11)\]

This value appears as a free parameter in the Standard Model. It is
measured and inserted. No Standard Model calculation derives it from the
theory's structure. Arthur Eddington famously argued on numerological
grounds that \(1/\alpha\) is exactly 137; Feynman described \(\alpha\)
as ``one of the greatest damn mysteries of physics.'' Prior
numerological attempts are cautionary context: Eddington's derivation
failed because he chose a structure to fit the number rather than
deriving the number from the structure.

The tholonic approach imposes a structural constraint before any
comparison with measurement: the expression must be built exclusively
from the five tholonic constants (\(\pi/4\), \(\phi\), \(e\),
\(\sqrt{2}\), \(\ln 2\)) and thologram structural integers (the axis
multipliers 2, 3, 5; the outer vertex values 1, 2, 4; and the role
weights 2, 3, 6), with all exponents being those fixed integers. No
fitting is permitted after the fact.

A systematic computational search over all such expressions identifies a
leading candidate:

\[\frac{1}{\alpha} \approx 4 \cdot e^5 \cdot (\ln 2)^4 = 137.035865 \quad \text{(error } {-0.98 \text{ ppm)}}\]

Since \((\sqrt{2})^4 = 2^2 = 4\) exactly, this is equivalently:

\[\frac{1}{\alpha} \approx e^5 \cdot (\sqrt{2})^4 \cdot (\ln 2)^4\]

a product of three tholonic constants with integer exponents. Its
tholonic structure is:

\[\frac{1}{\alpha} \approx C_{\text{outer}} \cdot e^{D_{\text{axis}}} \cdot (\ln 2)^{C_{\text{outer}}}\]

where \(C_{\text{outer}} = 4 = 2^2\) is the C outer vertex value,
\(D_{\text{axis}} = 5\) is the Definition axis multiplier, and the same
\(C_{\text{outer}} = 4\) reappears as the exponent of \(\ln 2\). All
three structural numbers are fixed thologram values. No alternating
series with tholonic seeds, no two-constant product with integer
exponents up to \(\pm 8\), and no four-constant product within the
searched range came within 200 ppm of \(1/\alpha\). This candidate is
the unique expression in the three-tholonic-constant space within 1 ppm
of the measured value.

The structural reading of the formula encodes the thologram's D-C
relationship: the D axis multiplier (5) governs the exponential rate via
\(e\), while the C outer vertex value (4) appears twice, as both the
scale factor and the power of \(\ln 2\). The C role is the accumulating,
integrating component; \(\ln 2\) is the tholonic constant associated
with logarithmic (integrating) convergence (Paper 1, Class B branch).
The D role is the bounding component; \(e\) raised to the D-axis power
sets the overall magnitude. Whether this structural reading can be
derived from role axioms rather than observed numerically is the central
open theoretical question.

\textbf{Status of the claim.} The leading candidate is 0.98 ppm from the
measured value. It is not exact at CODATA precision (\(0.15\) ppb). It
is the unique candidate in the constrained search space. A tholonic
derivation of \(\alpha\) should either produce an expression in the
neighborhood of
\(C_{\text{outer}} \cdot e^{D_{\text{axis}}} \cdot (\ln 2)^{C_{\text{outer}}}\)
or explain why that expression coincides to 1 ppm with the measured
value while being structurally motivated.

\textbf{Secondary candidate.} A secondary expression
\(4\pi^3 + \pi^2 + \pi = 137.036304\) (error \(+2.22\) ppm) can be read
as \(C \cdot \pi^{D+N} + N \cdot \pi^D + N \cdot \pi^N\) with outer
vertex values \(N = 1\), \(D = 2\), \(C = 4\). It is 2.3 times less
accurate than the leading candidate and the D vertex is absent from the
prefactors, which requires structural justification. It is listed here
for completeness.

\subsubsection{4.2 The proton/electron mass
ratio}\label{the-protonelectron-mass-ratio}

The ratio \(m_p / m_e \approx 1836.15267\) is not derived in the
Standard Model. The empirical approximation \(6\pi^5 \approx 1836.118\)
is known but unexplained. In the tholonic assignment for the hydrogen
atom (Paper 8, Section 4), the proton occupies the D role, the electron
occupies the C role, and the stable orbital (N state) emerges from their
balance. The role weights in the 2-3-6 pattern assign N = 6, D = 2, C =
3.

The approximation \(6\pi^5 \approx 1836\) has a natural tholonic
reading: the N weight (6) multiplied by \(\pi\) (the tholonic constant
most closely associated with the N recursion limit, the \(\pi/4\) branch
from Paper 1) raised to the D-axis power (5). Written with thologram
structural integers:

\[\frac{m_p}{m_e} \approx N_{\text{weight}} \cdot \pi^{D_{\text{axis}}} = 6 \cdot \pi^5 = 1836.118 \quad \text{(error } {-0.019\%)}\]

This is a 190 ppm approximation, worse than the fine structure constant
candidate but structurally analogous: both use the D-axis multiplier (5)
as an exponent and thologram role integers as prefactors. The
coincidence is suggestive. A derivation from role axioms would need to
show why the D-axis integer appears as the exponent of \(\pi\) in the
ratio of the D-role mass to the C-role mass. This is an open derivation.

\subsubsection{4.3 The Koide formula: D:C weight ratio and lepton
masses}\label{the-koide-formula-dc-weight-ratio-and-lepton-masses}

The Koide formula (discovered empirically in 1982) states that the three
charged lepton masses satisfy:

\[Q = \frac{m_e + m_\mu + m_\tau}{(\sqrt{m_e} + \sqrt{m_\mu} + \sqrt{m_\tau})^2} = \frac{2}{3}\]

This equality holds to current measurement precision
(\(Q_{\text{lepton}} = 0.66659\) using PDG 2024 values). The Standard
Model has no derivation of this value; it is one of the most prominent
unexplained numerical patterns in particle physics.

The tholonic 2-3-6 role weights assign D = 2 and C = 3, giving a D:C
weight ratio of exactly \(2/3\). If the three lepton generations occupy
the N, D, C roles in the lepton tholon, the Koide formula's \(Q = 2/3\)
is numerically identical to the tholonic D:C weight ratio. The value is
not imported from measurement; it is the structural weight ratio of the
tholonic roles.

\textbf{The quark triplet constraint.} Paper 8 (Section 8.1) reports the
Koide \(Q\) values for the quark triplets using PDG 2024 masses:
\[Q_{uct} \approx 0.849 \quad Q_{dsb} \approx 0.731\]

Both values are substantially above \(2/3\). This is a partial
falsification of the universal Koide claim. The tholonic identification
\(Q = D/C\) currently requires restriction to the lepton sector, with a
structural justification for why quarks do not satisfy the same
relation. Three candidate explanations are: (1) the Koide relation is a
property of the lepton sector specifically, where color charge is
absent; (2) quark mass running scheme dependence accounts for the
deviation at a unified renormalization scale; (3) the structural
derivation of the Koide form, once produced, will identify which
triplets are expected to satisfy \(Q = 2/3\). Until a structural
derivation of the Koide quadratic mean form exists, this remains a
suggestive correspondence rather than a full structural prediction.

\subsubsection{4.4 The Tsirelson bound: summary of the Paper 8
derivation}\label{the-tsirelson-bound-summary-of-the-paper-8-derivation}

Paper 8 (Section 3.4) derives the Tsirelson bound \(2\sqrt{2}\) from
tholonic first principles at two independent levels. The derivation is
summarized here because it constitutes a completed structural derivation
of a previously unexplained quantum bound.

\textbf{From the real-virtual tholon pair.} Paper 1 proves that the
tholonic recursive system converges to \(\sqrt{2}\) as one of its five
classical limits. A complete tholon always comprises a real N-D-C
configuration and its virtual complement. Both are full recursive
triadic systems and both independently converge to \(\sqrt{2}\). Their
composite value is \(\sqrt{2} + \sqrt{2} = 2\sqrt{2}\). An entangled
particle pair is, in tholonic terms, a complete real-virtual tholon
pair; the maximum achievable CHSH correlation is bounded by this
composite convergence.

\textbf{From pure angle geometry.} The tholonic D-C balance angle is
\(\pi/4\) (the optimal angle where \(\sin\theta = \cos\theta\)), which
is also the \(\pi\)-branch seed angle from Paper 1. At this angle, the
maximum CHSH value is:
\[\text{CHSH}_{\max} = 4\cos(\pi/4) = \frac{4}{\sqrt{2}} = 2\sqrt{2}\]

The CHSH sign pattern (\(+1, +1, +1, -1\)) is derived from three
tholonic axioms, the last of which (the D-C balance sign rule) is proved
from Paper 3, Lemma 4.2, without reference to measurement scenarios or
quantum mechanics. This is the first structural derivation of the CHSH
operator form and Tsirelson value from non-quantum-mechanical first
principles.

\begin{center}\rule{0.5\linewidth}{0.5pt}\end{center}

\subsection{5. Parameter Comparison: Standard Model vs Tholonic
Model}\label{parameter-comparison-standard-model-vs-tholonic-model}

\subsubsection{5.1 Comparison table}\label{comparison-table}

The table below lists the Standard Model's major free parameters and
structural posits alongside the tholonic model's current coverage.
``Derived'' means the tholonic model generates the value from structural
geometry without a free parameter. ``Structural'' means the tholonic
model provides a principled reason for the observed fact without
importing it from measurement. ``Open'' means the tholonic model has a
candidate or argument but lacks a completed derivation.

\begin{longtable}[]{@{}
  >{\raggedright\arraybackslash}p{(\linewidth - 6\tabcolsep) * \real{0.2500}}
  >{\raggedright\arraybackslash}p{(\linewidth - 6\tabcolsep) * \real{0.2500}}
  >{\raggedright\arraybackslash}p{(\linewidth - 6\tabcolsep) * \real{0.2500}}
  >{\raggedright\arraybackslash}p{(\linewidth - 6\tabcolsep) * \real{0.2500}}@{}}
\toprule\noalign{}
\begin{minipage}[b]{\linewidth}\raggedright
SM Parameter or Structural Fact
\end{minipage} & \begin{minipage}[b]{\linewidth}\raggedright
SM Status
\end{minipage} & \begin{minipage}[b]{\linewidth}\raggedright
Tholonic Status
\end{minipage} & \begin{minipage}[b]{\linewidth}\raggedright
Paper
\end{minipage} \\
\midrule\noalign{}
\endhead
\bottomrule\noalign{}
\endlastfoot
Three fermion generations & Observed; no SM explanation & Structural:
three roles are the minimum for stable recursion & 3, 8 \\
Three color charges & Postulated (SU(3) axiom) & Structural: three N-D-C
roles require three-fold balance & 8 \\
Three spatial dimensions & Observed; no SM explanation & Structural:
tetrahedron is the minimum volumetric closure & 8 \\
Quark confinement & Derived from QCD running coupling & Structural:
incomplete tholon is unstable & 8 \\
Koide \(Q = 2/3\) (lepton masses) & Measured; no SM derivation & Derived
from D:C weight ratio (lepton sector) & 8 \\
CHSH Tsirelson bound \(2\sqrt{2}\) & Derived from Hilbert space algebra
& Structural derivation from tholonic first principles & 8 \\
Fine structure constant \(1/\alpha \approx 137.036\) & Measured; no SM
derivation & Open: leading candidate \(4 \cdot e^5 \cdot (\ln 2)^4\),
error \(-0.98\) ppm & this paper \\
Proton/electron mass ratio \(\approx 1836\) & Measured; no SM derivation
& Open: candidate \(6\pi^5 \approx 1836.12\), error \(-190\) ppm & this
paper \\
Number of SM quantum fields (17) & Postulated independently & Open:
reduction if color, generation, composition structure all follow from
one tholonic recursion & 3 \\
Quark and lepton masses (12 parameters) & Measured; no SM derivation &
Koide lepton sector derived; quark sector constrained by \(Q_{uct}\),
\(Q_{dsb}\) values & 8 \\
CKM mixing angles (4 parameters) & Measured; no SM derivation & Not yet
addressed & -- \\
Gauge coupling constants (3 parameters) & Measured; no SM derivation &
\(\alpha\) candidate above; weak and strong couplings not yet addressed
& this paper \\
Higgs mass and quartic coupling (2 parameters) & Measured; no SM
derivation & Not yet addressed & -- \\
Strong CP phase (near zero; unexplained) & Measured near zero; no SM
derivation & Not yet addressed & -- \\
\end{longtable}

\subsubsection{5.2 What remains to be
derived}\label{what-remains-to-be-derived}

The tholonic model's current position is strongest on structural facts
(three generations, three colors, three dimensions, confinement, Koide
\(2/3\), Tsirelson bound) and weakest on quantitative parameters where
the Standard Model's fine-grained machinery (renormalization group,
gauge structure, Higgs mechanism) is most powerful. The CKM mixing
angles, Higgs parameters, and strong CP phase have no current tholonic
address. The weak and strong coupling constants beyond \(\alpha\) are
not yet addressed quantitatively.

This is an honest accounting. A framework that addresses five of the
Standard Model's structural silences and has a 1 ppm candidate for the
most prominent free parameter is in a different position than a
framework that merely claims consistency. But it is also not a complete
alternative. It is a candidate alternative, with a specified list of
open derivations that would, if completed, move it into the position of
a genuine competitor.

\begin{center}\rule{0.5\linewidth}{0.5pt}\end{center}

\subsection{6. Falsifiability}\label{falsifiability}

The tholonic model is falsifiable. The following conditions are
organized by type.

\subsubsection{6.1 Strong falsification
conditions}\label{strong-falsification-conditions}

\textbf{Fine structure constant.} If a derivation from thologram
geometry produces a value that differs from the measured
\(1/137.035999206\), that specific derivation is falsified. If an
exhaustive search over all tholonic-constant expressions with structural
integers as exponents finds no candidate within, say, 10 ppm, the claim
that \(\alpha\) is a tholonic balance ratio is falsified at that scope.

\textbf{Fourth fermion generation.} If a fourth-generation charged
fermion is discovered with mass accessible to current or near-future
colliders, the structural prediction of exactly three generations is
falsified.

\textbf{Isolated free quark.} If a stable, isolated free quark or
color-nonsinglet hadron is confirmed experimentally, the
incomplete-tholon instability argument for confinement is falsified.

\textbf{Koide formula in lepton sector.} If new precision measurements
of the charged lepton masses show \(Q_{\text{lepton}}\) deviating from
\(2/3\) by more than measurement uncertainty, the Koide correspondence
is falsified at the level of precision measurement.

\subsubsection{6.2 Structural constraints}\label{structural-constraints}

\textbf{N-D-C non-assignability.} If a stable, self-sustaining,
recursive physical system is found that cannot be consistently assigned
an N-D-C structure at any scale, the universality claim of the tholonic
framework requires explicit revision. The framework would need to
specify which systems it applies to and why.

\textbf{Koide quark sector.} The quark triplet \(Q\) values
(\(Q_{uct} \approx 0.849\), \(Q_{dsb} \approx 0.731\)) already
constitute a partial constraint. The tholonic model must either (a)
restrict the Koide D:C correspondence to leptons with structural
justification, or (b) demonstrate that a renormalization-group or
color-corrected quark mass evaluation recovers \(Q \approx 2/3\).
Failure of both options weakens the Koide claim structurally.

\subsubsection{6.3 What would constitute
confirmation}\label{what-would-constitute-confirmation}

A completed derivation of \(1/\alpha\) from thologram geometry,
producing a value within the CODATA measurement precision of \(0.15\)
ppb, would constitute the single strongest available confirmation. It
would demonstrate that a dimensionless physical constant is determined
by the geometry of the tholonic recursion structure, which is precisely
the criterion Einstein identified for a genuinely unified field theory.

A completed derivation of the proton/electron mass ratio from N-D-C role
weights would constitute a second confirmation of the same type.

Discovery that the quark sector Koide \(Q\) values approach \(2/3\) at a
unified renormalization scale would substantially extend the
lepton-sector Koide correspondence to the complete fermion spectrum.

\begin{center}\rule{0.5\linewidth}{0.5pt}\end{center}

\subsection{7. The Path Forward}\label{the-path-forward}

\subsubsection{7.1 The fine structure constant derivation as the
priority
target}\label{the-fine-structure-constant-derivation-as-the-priority-target}

The fine structure constant derivation is the highest priority because
it is the most testable and most consequential open item. The leading
candidate \(4 \cdot e^5 \cdot (\ln 2)^4\) is at 0.98 ppm. The gap
between a numerical coincidence and a structural derivation is the gap
between observing that the expression is close and proving that the
structure of the thologram constrains \(\alpha\) to take that value.

The direction for a structural derivation is suggested by the formula's
own structure: \(D_{\text{axis}} = 5\) appears as the exponent of \(e\)
(the exponential growth constant), while \(C_{\text{outer}} = 4\)
appears both as the prefactor and as the exponent of \(\ln 2\) (the
logarithmic integration constant). The D-C duality in the exponent
structure mirrors the D-C functional duality of the tholonic roles. A
derivation from role axioms would need to show why the electromagnetic
coupling strength, which governs how strongly charged fields interact,
is determined by the ratio of the thologram's D-axis exponential scaling
to its C-role logarithmic convergence.

One approach is dimensional: the fine structure constant is
dimensionless, a ratio. Dimensionless quantities in physics characterize
the balance between competing tendencies. The tholonic model is
precisely a theory of how competing tendencies (D and C) reach a stable
balance (N). An electromagnetic coupling of \(\alpha \approx 1/137\)
says that the electromagnetic force is very weak compared to the
fundamental unit of quantum charge. In tholonic terms, this is a
statement about how far from the D-C balance point the electromagnetic
interaction sits in the thologram. The derivation challenge is to make
this structural argument quantitative.

\subsubsection{7.2 Required theoretical
work}\label{required-theoretical-work}

In order of priority:

\begin{enumerate}
\def\labelenumi{\arabic{enumi}.}
\tightlist
\item
  \textbf{Fine structure constant structural derivation.} Move from the
  0.98 ppm numerical candidate to a role-axiomatic derivation that
  produces the expression from D-C balance geometry without fitting.
\item
  \textbf{Koide formula structural derivation.} Produce a derivation
  showing that the quadratic mean form of the Koide formula follows from
  the tholonic role weights, and identify which physical systems are
  expected to satisfy it.
\item
  \textbf{Proton/electron mass ratio derivation.} Show why
  \(N_{\text{weight}} \cdot \pi^{D_{\text{axis}}}\) gives the ratio of
  D-role to C-role masses in the hydrogen tholon.
\item
  \textbf{Three spatial dimensions formal proof.} Prove formally that
  stable recursive N-D-C structure with volumetric closure requires
  exactly three spatial dimensions.
\item
  \textbf{Quark mass Koide deviation.} Determine whether a unified
  renormalization scale or a color-charge correction to the tholonic
  mapping recovers \(Q \approx 2/3\) for quark triplets, or formally
  restrict the Koide correspondence to the lepton sector.
\item
  \textbf{CKM mixing angle address.} Identify whether the quark mixing
  hierarchy (Cabibbo angle, b-quark suppression) has a tholonic
  structural reading.
\end{enumerate}

\subsubsection{7.3 Scope of the claim}\label{scope-of-the-claim}

This paper claims that the tholonic model is a \emph{candidate}
alternative to the Standard Model, not a replacement. The distinction
matters. A candidate alternative is a framework that (a) is internally
consistent, (b) addresses at least some of the target framework's open
problems, (c) makes falsifiable predictions that the target framework
does not make, and (d) has a specified path toward the open derivations
that would, if completed, constitute the strongest available
confirmation.

The tholonic model currently satisfies all four conditions. It does not
currently satisfy the stronger criterion: it does not yet derive enough
Standard Model parameters with sufficient precision to constitute a
predictive replacement. That criterion requires completing at least the
fine structure constant derivation. Until that is done, the correct
characterization is: structurally better in several respects,
quantitatively competitive in one, and with a clear agenda for the work
required to make the quantitative case conclusive.

\begin{center}\rule{0.5\linewidth}{0.5pt}\end{center}

\subsection{8. Implications}\label{implications}

The Standard Model's 19 free parameters and 17 independently postulated
fields are not merely aesthetic deficiencies. They represent genuine
explanatory gaps: properties of the physical world that the most
successful theory in physics treats as given rather than derived. The
historical analogy is instructive. Kepler's model of planetary orbits
was more predictively accurate than Copernicus's circles; Newton's
mechanics explained why Kepler's laws hold. The Standard Model is in
Kepler's position: extraordinarily accurate, but explaining rather than
deriving the structure it describes.

The tholonic model proposes that the structure underlying the Standard
Model is the N-D-C tholonic recursion. The proposal is speculative. But
it is internally consistent, makes falsifiable predictions, and
addresses several of the Standard Model's structural silences without
multiplying free parameters. Its five structural correspondences (three
generations, three colors, three spatial dimensions, confinement, Koide
\(2/3\)) come at the cost of one additional axiom: the N-D-C triad as
the universal unit of stable recursive organization. Whether that
additional axiom is worth the explanatory gain is an empirical question,
and its answer depends on whether the quantitative derivations in
Section 7.2 succeed.

Tesla and Einstein both argued, from different positions, that the field
is primary and that what we call particles are stable field
configurations. The tholonic model provides a specific structural
grammar for that claim: the A\(\&\)I field is the single underlying
field, the tholonic recursion is the organizing principle, and the
Standard Model's diverse particle spectrum is the thologram's stable
excitation structure. The unified field that Einstein pursued for thirty
years is proposed to be this structure. Whether it is that structure is
the question the derivations in Section 7.2 will decide.

\begin{center}\rule{0.5\linewidth}{0.5pt}\end{center}

\subsection{References and Source
Materials}\label{references-and-source-materials}

\begin{itemize}
\tightlist
\item
  \textbf{Tholonia: The Mechanics of Existential Awareness} by Duncan
  Stroud, especially Chapters 11 through 15.
\item
  \textbf{Paper 1 in this series:} \emph{Emergence of Classical
  Constants from a Minimal Recursive Triadic Framework.} Derivation of
  \(\pi/4\), \(\phi\), \(e\), \(\sqrt{2}\), \(\ln 2\) from tholonic
  recursion.
\item
  \textbf{Paper 3 in this series:} \emph{Minimal Recursive Triadic
  Framework.} Formal proof that three functional roles are the minimum
  for non-trivial convergent recursion.
\item
  \textbf{Paper 6 in this series:} \emph{The Qualitative Nature of One,
  Two, and Three: Structural Role Assignment in Minimal Recursive
  Systems.} Formal basis for the N-D-C integer role assignment.
\item
  \textbf{Paper 7 in this series:} \emph{Engineering Toward N: Cambridge
  Semantics as Empirical Validation of the Tholonic N-D-C Framework.}
  Independent engineering convergence on the N-D-C structure.
\item
  \textbf{Paper 8 in this series:} \emph{The Atom as a Measurable
  Tholon: From Einstein's Incomplete Electron to a Field-First
  Ontology.} Atomic correspondences and Standard Model comparison that
  this paper extends.
\item
  \textbf{Koide, Y.} (1982). ``A fermion-boson composite model of quarks
  and leptons.'' \emph{Physics Letters B} 120, 161--165.
\item
  \textbf{Cirel'son, B.S.} (1980). ``Quantum generalizations of Bell's
  inequality.'' \emph{Letters in Mathematical Physics} 4, 93--100.
\item
  \textbf{Feynman, R.P.} (1985). \emph{QED: The Strange Theory of Light
  and Matter.} Princeton University Press. (Source of Feynman's
  description of \(\alpha\) as ``one of the greatest damn mysteries of
  physics.'')
\item
  \textbf{Einstein, A.} (1949). ``Autobiographical Notes.'' In
  \emph{Albert Einstein: Philosopher-Scientist}, ed.~P.A. Schilpp. Open
  Court.
\item
  \textbf{Particle Data Group.} (2024). \emph{Review of Particle
  Physics.} Progress of Theoretical and Experimental Physics 2024,
  083C01.
\item
  \textbf{CODATA.} (2022). \emph{CODATA Internationally Recommended 2022
  Values of the Fundamental Physical Constants.} NIST.
\item
  \textbf{Weinberg, S.} (1995). \emph{The Quantum Theory of Fields},
  Vol. 1. Cambridge University Press.
\item
  \textbf{Gross, D.J., Politzer, H.D., Wilczek, F.} (1973).
  ``Ultraviolet Behavior of Non-Abelian Gauge Theories.'' \emph{Physical
  Review Letters} 30, 1343--1346. (Asymptotic freedom, QCD basis for
  confinement.)
\item
  \textbf{ALEPH, DELPHI, L3, OPAL Collaborations.} (2006). ``Precision
  electroweak measurements on the Z resonance.'' \emph{Physics Reports}
  427, 257--454. (LEP neutrino species constraint
  \(N_\nu = 2.9840 \pm 0.0082\).)
\end{itemize}

\end{document}
